\section{Introduction}\label{sec:intro}

Networked extended reality (XR) changed shape after 2023. The standalone headset became the main client:
it tracks the user, renders what it can, and streams the rest from an edge or cloud renderer. Split
rendering moved from cables to Wi-Fi~6 and 7 and to 5G-Advanced, where 3GPP Releases 18 and 19 added the
first radio features designed around XR traffic \cite{W7164211432}. The content changed too.
Besides rendered frames and 360-degree video, systems now stream point clouds, avatars sent as keypoints
\cite{W4403981191}, and 3D Gaussian splats \cite{W4416477288}. The metaverse push added social sessions in
which several people share one scene and judge each other's delay.

These changes couple networking and quality of experience (QoE) more tightly than before. A frame is
useful only if it arrives before its display deadline, so the network decides whether an image appears,
not only how sharp it is. The headset hides head-rotation delay by reprojection, but it cannot hide a lost
frame, a frozen peer or a stalled scene. The coupling now also runs the other way. The radio can learn
which packets form one frame and how much each frame matters \cite{W4410582180}, and an application can be
told what rate the cell can carry \cite{W7161282618}. \Cref{fig:loop} draws the resulting control loop.
Every mechanism in this review closes some part of it, and the question for each is whether the part it
closes is the one users perceive.

\begin{figure}[t]
\centering
\begin{adjustbox}{max width=\linewidth}\begin{tikzpicture}[font=\sffamily\footnotesize,
  blk/.style={draw=black!60, fill=black!3, rounded corners=2pt, align=center, text width=2.65cm,
    minimum height=17mm, inner sep=2pt, line width=0.6pt},
  usr/.style={blk, draw=gapcolor!80!black, fill=gapcolor!5},
  data/.style={-{Stealth[length=1.8mm]}, line width=0.7pt, draw=black!70},
  app/.style={-{Stealth[length=1.8mm]}, line width=0.8pt, draw=encode!85!black},
  net/.style={-{Stealth[length=1.8mm]}, line width=0.8pt, draw=provision!85!black, densely dashed},
  aware/.style={-{Stealth[length=1.8mm]}, line width=0.8pt, draw=compute!75!black, dash dot},
  lab/.style={font=\sffamily\scriptsize, align=center, text=black, inner sep=1.5pt}]
\node[usr] (u) at (0,0)     {\textbf{User}\\interaction delay, stalls, quality, sickness (\S\ref{sec:qoe})};
\node[blk] (h) at (3.8,0)   {\textbf{Headset}\\decode, reprojection; pose, gaze, frame statistics (\S\ref{sec:measure})};
\node[blk] (a) at (7.6,0)   {\textbf{Radio access}\\PDU sets, power saving, \mbox{Wi-Fi~7}, mmWave (\S\ref{sec:radio})};
\node[blk] (t) at (11.4,0)  {\textbf{Transport}\\L4S/ECN, QUIC, MoQ, \mbox{WebRTC}, FEC (\S\ref{sec:transport})};
\node[blk] (r) at (15.2,0)  {\textbf{Edge/cloud renderer}\\encoder, bitrate, foveation, splats (\S\ref{sec:transport})};
% data plane: downlink above, uplink below
\foreach \x/\y in {r/t,t/a,a/h} \draw[data] ([yshift=4mm]\x.west) -- ([yshift=4mm]\y.east);
\foreach \x/\y in {h/a,a/t,t/r} \draw[data] ([yshift=-4mm]\x.east) -- ([yshift=-4mm]\y.west);
\draw[data, {Stealth[length=1.8mm]}-{Stealth[length=1.8mm]}] (u.east) -- (h.west);
% control loops
\draw[app] (h.north) -- ++(0,1.25) -| node[lab, pos=0.25, above] {QoE estimate, gaze, loss and delay reports
  $\rightarrow$ bitrate, foveation, what to render} (r.north);
\draw[net] ([xshift=4mm]a.north) -- ++(0,0.7) -| node[lab, pos=0.27, above] {rate recommendation, ECN marks} ([xshift=-5mm]r.north);
\draw[aware] ([xshift=-5mm]r.south) -- ++(0,-0.65) -| node[lab, pos=0.25, below] {PDU-set boundaries, importance, burst timing $\rightarrow$ frame-aware scheduling} (a.south);
\draw[app] (u.south) -- ++(0,-0.65) -| node[lab, pos=0.27, below] {ratings $\rightarrow$ QoE models} (h.south);
% legend
\node[lab, anchor=west, align=left, font=\sffamily\footnotesize] at (-1.35,-2.35) {%
  \tikz[baseline=-0.6ex]\draw[data] (0,0)--(0.6,0); data (frames down; pose, gaze, input up)\quad
  \tikz[baseline=-0.6ex]\draw[app] (0,0)--(0.6,0); application feedback\quad
  \tikz[baseline=-0.6ex]\draw[net] (0,0)--(0.6,0); network $\rightarrow$ application\quad
  \tikz[baseline=-0.6ex]\draw[aware] (0,0)--(0.6,0); application $\rightarrow$ network};
\end{tikzpicture}\end{adjustbox}
\caption{The QoE control loop of a networked XR session and the sections that discuss each part. The
application adapts to its own measurements, the network recommends rates or marks congestion, and the
application tells the network which packets form a frame. Only the path from users' ratings to QoE models
rests on user studies.}
\label{fig:loop}
\end{figure}

This review covers work published from 2023 to 2026 and follows the loop from the user outwards:
what users perceive (\cref{sec:qoe}), what measurements of real systems show (\cref{sec:measure}), the
radio access (\cref{sec:radio}), and transport and delivery (\cref{sec:transport}). \Cref{tab:families}
classifies the families of mechanisms, and \cref{sec:outlook} states open problems. Throughout, we
separate what is \emph{established} (measured on real systems or with users), what is \emph{contested},
and what is \emph{assumed}.
